\chapter*{整理者前言}
\addcontentsline{toc}{chapter}{整理者前言}

我很早就想把谢惠民等编著的《数学分析习题课讲义》完整整理成一份方便查阅的电子版。这个念头拖了很久，直到现在才真正做出来。v1.0 合订版最后长到了 1512 页，我私下把它叫作“谢惠民的方便面”。希望需要查一道题、补一个证明或对照一个解答时，拿来就能用。

早些年，网上比较容易遇到的习题解答主要有两套。一套通常署名“顺数人”，只覆盖上册，但曾经免费分享；另一套常被称作“甲乙版”，覆盖范围更广。我当时能接触到的公开资料里，甲乙版流传出来的多是部分内容，完整版本据当时的说明主要以带购买者身份信息水印的纸质形式提供，获取门槛较高。这些材料都很有价值，也让我一直希望能有一份完整、可检索、可以持续修订的版本。

AI 工具成熟以后，这个计划终于变得可行。本项目先用 AI 辅助把上下册扫描内容逐章整理成结构化 LaTeX，正文和题面仍以原书扫描页为依据。随后让 AI 在不依赖流传答案的前提下逐题作答。初稿完成后，再用“顺数人”和“甲乙”两套解答逐题交叉复核。参考答案中更简洁、更自然或更有启发性的做法会被吸收到正式解答里；遇到彼此冲突的地方，则重新推导判断。这样做的目的，是让最终版本保留独立作答的完整过程，同时尽量吸收已有解答中真正值得留下的思路。

为了方便连续阅读，本版把第 1 至 26 章排成一部完整讲义，练习题、思考题和参考题的解答直接跟在原题之后。正文与题面尽量忠实于扫描本，数学记号和版式则重新统一。需要题图时优先从原书裁取，不用 AI 重画替代原图。

\section*{致谢}

首先感谢谢惠民、恽自求、易法槐、钱定边等原作者留下这套内容扎实、题目丰富的习题课讲义。

感谢“顺数人”曾经免费分享上册解答，也感谢“甲乙版”解答的整理者。两套材料为这次逐题复核提供了重要参照，其中不少证明和计算思路直接启发了本版的改写与另解。也感谢所有曾经在网上分享相关资料、讨论题目和留下勘误的人。

本项目大量使用 AI 工具完成转录、初次作答、交叉检查和排版整理。它让一项原本很难靠个人手工完成的工作真正落了地。即使经过多轮复核，一千多道题仍难免留下笔误、疏漏或可以继续改进的解法。若读者发现问题，欢迎留下题号、页码和推导，一起把这份版本继续修好。
